% Response to reviewers for manuscript jimaging-4584529.
% Locators verified against last_2026-10-17_10-17-58.tex, MDPI highlighted PDF (23 pages).
\documentclass[11pt,a4paper]{article}
\usepackage[margin=2.4cm]{geometry}
\usepackage{amsmath}
\usepackage{xcolor}
\usepackage{enumitem}
\usepackage{url}
\newenvironment{revcomment}{\par\medskip\noindent\begin{minipage}{\linewidth}\itshape}{\end{minipage}\par\smallskip}
\setlength{\parindent}{0pt}
\setlength{\parskip}{4pt}
\setlength{\emergencystretch}{3em}

\begin{document}

\begin{center}
{\Large\bfseries Response to Reviewers}\\[4pt]
{\large Refining Retinal Vessel Segmentation with Conditional Diffusion Models:\\ A Controlled Evaluation on LU-Net+RA Probability Maps}\\[4pt]
Journal of Imaging \textendash{} Manuscript ID: jimaging-4584529
\end{center}

Dear Editor,

We thank both reviewers for their detailed comments. In addressing them, we re-ran the study under a fully documented protocol, with repeated training, saved predictions and evaluation of refined and coarse predictions on the same test images. This changed the main result, and we have revised the manuscript accordingly. The main changes are:
\begin{enumerate}[leftmargin=*,itemsep=2pt]
\item \textbf{Controlled evaluation.} Refined and coarse LU-Net+RA predictions are now compared on the same test images, with five training seeds and five sampling seeds for each stochastic repeated-training configuration. Under this protocol the refiner did not improve mean Dice: 0.746 versus 0.811 on CHASE-DB1, 0.790 versus 0.792 on DRIVE and 0.666 versus 0.732 on HRF (Section~4.4, Table~6).
\item \textbf{Earlier results.} The checkpoint and predictions behind the single-run scores of the original submission (Dice 0.835, 0.835 and 0.793) are no longer available, and the re-implementation did not reproduce them. We no longer report those scores as results.
\item \textbf{Ablations and baselines.} We added component ablations and a spatial residual refiner (Sections~4.5 and~4.6).
\item \textbf{Paired uncertainty and OOF.} Section~4.4.1 now provides all per-image Dice means and paired changes (Tables~7--9) and paired 95\% intervals (Table~6). Section~4.4.2 reports the completed matched OOF comparison (Table~10), with its statistical and design limitations.
\item \textbf{Reporting.} We added a parameter audit, a layer-level denoiser specification, the full training and evaluation protocol, a matched 50- versus 100-update comparison and a complete-pipeline latency benchmark (Sections~3.3, 4.1--4.3, 4.8 and~4.9).
\item \textbf{Title.} Because the results no longer support the description ``a lightweight enhancement'', we kept the main title and replaced the subtitle with ``A Controlled Evaluation on LU-Net+RA Probability Maps''.
\end{enumerate}
Changes are highlighted in blue in the revised manuscript. Page and line numbers below refer to the highlighted version.

\section*{Reviewer 1}
\subsection*{Comment 1}
\begin{revcomment}The reported improvement over LU-Net+RA cannot be verified quantitatively. For CHASE-DB1, the proposed method reports Dice = 0.835, compared with 0.7946 for the published LU-Net+RA result, a nominal gain of 0.0404. However, these values come from different test protocols. The paper should report the actual Dice of LU-Net+RA on the same 8 CHASE-DB1 test images before claiming improvement.\end{revcomment}
\textbf{Response.} We agree that the published LU-Net+RA score cannot measure the refinement effect. We now evaluate the coarse LU-Net+RA predictions and the refined predictions on the same eight CHASE-DB1 test images (four subjects), with identical preprocessing, annotations, threshold and metric code, five training seeds and five sampling seeds. Mean Dice is $0.8105\pm0.0018$ for LU-Net+RA and $0.7461\pm0.0026$ after refinement, a difference of $-0.064$. The refiner therefore does not improve the baseline on CHASE-DB1. The value of 0.835 in the original submission came from a single run whose checkpoint is no longer available, and the re-implementation did not reproduce it; we no longer report it as a result. The published value of 0.7946 is kept only as literature context.

\textbf{Changes in the manuscript.} Section~4.4 (page~12, lines~418--432); Section~4.4.1 (pages~12--13, lines~434--453); Section~4.7 (page~16, lines~539--550); Table~6 (page~12); Table~13 (page~17).

\subsection*{Comment 2}
\begin{revcomment}The DRIVE comparison is similarly difficult to interpret. The reported Dice is 0.835 versus 0.7871 for published LU-Net+RA, a difference of 0.0479, but the sensitivity of the proposed model (0.825) is substantially higher than the published LU-Net+RA sensitivity (0.7421). Because the evaluation protocols are different, this apparent improvement should not be interpreted as an effect size without same-image testing.\end{revcomment}
\textbf{Response.} On the same 20 DRIVE test images, mean Dice is $0.7918\pm0.0076$ for LU-Net+RA and $0.7898\pm0.0083$ after refinement (difference $-0.002$), so the refiner does not improve DRIVE either. The sensitivity values quoted in the comment came from the original single-run results, which could not be reproduced and are no longer reported as results.

\textbf{Changes in the manuscript.} Section~4.4 (page~12, lines~418--432); Section~4.4.1 (pages~12--13, lines~434--453); Table~6 (page~12); Table~8 (page~13).

\subsection*{Comment 3}
\begin{revcomment}The HRF result is not consistently better. The proposed 20-step model obtains Dice = 0.762, which is actually 0.0136 lower than the published LU-Net+RA value of 0.7756. This weakens the general claim that the diffusion refinement consistently improves the baseline.\end{revcomment}
\textbf{Response.} The controlled evaluation confirms the reviewer's concern. On the same 12 HRF test images, mean Dice is $0.7321\pm0.0035$ for LU-Net+RA and $0.6657\pm0.0063$ after refinement. The refiner does not improve the baseline on any of the three datasets, and the abstract, discussion and conclusion no longer claim an improvement.

\textbf{Changes in the manuscript.} Abstract (page~1); Section~4.4 (page~12, lines~418--432); Section~5 (page~19, lines~599--630); Section~6 (page~20, lines~661--674); Table~6 (page~12).

\subsection*{Comment 4}
\begin{revcomment}The gain from increasing DDIM steps is relatively small after 50 steps. On CHASE-DB1, Dice increases from 0.835 at 50 steps to only 0.836 at 100 steps, while the reported refinement time doubles from approximately 100 ms to 200 ms. This suggests that 100 steps are computationally inefficient, but the paper should quantify whether the 0.001 difference is simply sampling variation.\end{revcomment}
\textbf{Response.} We ran a matched comparison on one trained CHASE-DB1 refiner: the eight test images were refined with 50 and with 100 DDIM updates from the same 20 initial-noise tensors. Mean Dice was 0.745476 and 0.745477 (difference $5.9\times10^{-7}$; 95\% bootstrap interval $[-9.4\times10^{-7},\,2.0\times10^{-6}]$), while mean refinement time rose from 116.98 to 229.64~ms. The interval includes zero and does not demonstrate a reliable Dice benefit from 100 updates at approximately twice the refinement cost; it does not establish equivalence. We retain 50 updates as the reference setting, without claiming it is optimal. The 0.835-to-0.836 difference in the original submission came from single runs that we no longer report.

\textbf{Changes in the manuscript.} Section~4.8 (page~17, lines~552--568); Section~4.9 (pages~17--18, lines~570--585).

\subsection*{Comment 5}
\begin{revcomment}The 10-to-50 step improvement comes at a substantial computational cost. Dice increases from 0.814 to 0.835 (+0.021), while runtime increases from approximately 20 ms to 100 ms (5$\times$). The paper should report a more formal accuracy--latency analysis rather than describing this only qualitatively.\end{revcomment}
\textbf{Response.} We replaced the approximate timings with measurements. A complete-pipeline benchmark gives 4.51~ms for LU-Net+RA, 117.95~ms for refinement, 119.02~ms for the complete model and 140.50~ms from a decoded RGB image to a binary mask in CPU memory, so the complete model takes about 26 times as long per image as the backbone. The matched comparison in our response to Comment~4 shows that doubling the update count approximately doubles refinement time with negligible observed Dice change. The 10-to-50-update figures in the comment came from the original single-run results, and update counts below 50 were not evaluated in the controlled protocol, so we do not report a full accuracy--latency curve.

\textbf{Changes in the manuscript.} Section~4.8 (page~17, lines~552--568); Section~4.9 (pages~17--18, lines~570--585); Table~14 (page~17).

\subsection*{Comment 6}
\begin{revcomment}The effect of DDIM steps is not monotonic across all metrics. On DRIVE, increasing the steps from 20 to 50 raises Dice from 0.795 to 0.835 (+0.040), but sensitivity decreases from 0.835 to 0.825 (-0.010). This indicates that the higher Dice does not correspond to uniformly better vessel detection and warrants analysis of FP/FN changes.\end{revcomment}
\textbf{Response.} We agree that higher Dice need not mean uniformly better vessel detection. The 20-to-50-update sensitivity change quoted in the comment came from the original single-run results, which could not be reproduced. The 20-update setting was not evaluated in the controlled protocol, so we cannot report paired FP/FN changes between 20 and 50 updates; the matched comparison of 50 and 100 updates shows negligible observed Dice change, with no demonstrated reliable benefit from 100 updates.

\textbf{Changes in the manuscript.} Section~4.4 (page~12, lines~418--432); Section~4.8 (page~17, lines~552--568).

\subsection*{Comment 7}
\begin{revcomment}The test sets are very small for drawing strong dataset-level conclusions. CHASE-DB1 has only 8 test images, while HRF has 12. With such small samples, a Dice change of 0.01--0.02 may not be statistically meaningful. Per-image results and 95\% confidence intervals should be reported.\end{revcomment}
\textbf{Response.} We now report the complete paired coarse-versus-diffusion analysis for the same global/frozen reference as Table~6. Tables~7--9 provide all 40 per-image Dice means and paired changes. Table~6 includes exploratory paired crossed-bootstrap 95\% intervals. The refined-minus-coarse differences and exploratory 95\% intervals are:
\begin{itemize}[leftmargin=1.5em,itemsep=2pt]
\item CHASE-DB1: $-0.064368245$; $[-0.078706654,-0.050060076]$.
\item DRIVE: $-0.002057107$; $[-0.003024955,-0.001308524]$.
\item HRF: $-0.066395394$; $[-0.074121150,-0.058866677]$.
\end{itemize}
Sampling realisations are averaged within image and training seed. The bootstrap uses 50,000 replicates and seed 20261017, resampling training seeds and test units while retaining both CHASE eyes within subjects and HRF pathology strata. All three intervals are below zero in this tested configuration. With only four CHASE subjects and five training pairs, and unverified patient independence for DRIVE and HRF, these intervals remain exploratory rather than confirmatory or simultaneous. They do not establish population-level robustness.

\textbf{Changes in the manuscript.} Section~4.3 (pages~11--12, lines~395--416); Section~4.4.1 (pages~12--13, lines~434--453); Table~6 (page~12); Table~7 (page~13); Table~8 (page~13); Table~9 (page~14); Section~5.1 (pages~19--20, lines~632--659).

\subsection*{Comment 8}
\begin{revcomment}The parameter-efficiency claim needs a more informative comparison. The proposed refinement module adds 502,401 parameters, but the manuscript does not provide a directly comparable parameter count for the LU-Net+RA baseline using the same counting convention. Therefore, it is difficult to judge how lightweight the complete system actually is.\end{revcomment}
\textbf{Response.} We audited all parameters with one convention: unique registered weights and biases, including frozen modules, excluding buffers and optimizer state. LU-Net+RA has 1,946,961 parameters. The executable refiner adds 527,233 (total 2,474,194; $+27.08\%$), of which 405,824 are updated during refinement training. The archived architecture adds 535,425 (total 2,482,386; $+27.50\%$); the 502,401 reported originally omitted a separate 33,024-parameter conditioning network. Parameter counts do not measure inference cost, which we now report separately.

\textbf{Changes in the manuscript.} Section~3.3 (page~5, lines~180--193); Table~1 (page~6); Table~2 (page~6); Section~1 (pages~1--2, lines~20--78).

\subsection*{Comment 9}
\begin{revcomment}The 128-dimensional global image embedding may be too restrictive for thin-vessel refinement. The diffusion model receives only a 128-dimensional global vector from the RGB image rather than spatial image features. Since the main motivation is recovery of local thin vessels and boundaries, an ablation comparing global conditioning with spatial conditioning would be useful.\end{revcomment}
\textbf{Response.} We added this ablation, with five independently trained refiners per configuration and dataset. Replacing the 128-dimensional global vector with spatial RGB features and spatial FiLM changed mean Dice by $+0.0002$ on CHASE-DB1, less than $0.0001$ on DRIVE and $+0.0001$ on HRF, and removing the image condition entirely had a similarly negligible effect. In this configuration, the tested conditioning changes have little effect; the frozen mask encoding is a plausible contributor to the loss, but these comparisons do not establish it as the sole cause.

\textbf{Changes in the manuscript.} Section~4.5 (pages~15--16, lines~481--514); Table~11 (page~15); Section~5 (page~19, lines~599--630).

\subsection*{Comment 10}
\begin{revcomment}The latent representation is not actually compressed in terms of scalar count. The input mask contains 262,144 values, whereas the encoded representation contains 1,048,576 values (4$\times$ more). Therefore, calling this a lightweight latent-space representation is misleading; it is spatially downsampled but has a larger total feature count.\end{revcomment}
\textbf{Response.} We agree and changed the terminology in the text and in Figure~1. The encoder maps the $1\times512\times512$ mask (262,144 values) to a $64\times128\times128$ grid (1,048,576 values): it reduces each spatial dimension fourfold but increases the scalar count fourfold. The manuscript now calls it a spatially downsampled multi-channel representation and makes no compression claim.

\textbf{Changes in the manuscript.} Section~1 (pages~1--2, lines~20--78); Section~3.3.1 (pages~5--7, lines~195--228); Figure~1 (page~3).

\subsection*{Comment 11}
\begin{revcomment}The contribution of the Dice term is not demonstrated. The diffusion objective uses $\lambda_{\mathrm{Dice}}$ = 0.5, but there is no ablation with $\lambda_{\mathrm{Dice}}$ = 0 or alternative weights. Since this term directly introduces segmentation supervision, its contribution should be separated from the diffusion denoising objective.\end{revcomment}
\textbf{Response.} We did not run a measured ablation of the Dice weight, so the contribution of the Dice term cannot be separated in this revision. The manuscript now states that $\lambda_{\mathrm{dice}}=0.5$ was fixed in advance and not tuned, and lists the missing ablation as a limitation.

\textbf{Changes in the manuscript.} Section~3.3.5 (pages~8--9, lines~285--317); Section~5.1 (pages~19--20, lines~632--659).

\subsection*{Comment 12}
\begin{revcomment}The claimed inference throughput is incomplete. The reported 10 images/s at 50 transitions excludes the LU-Net+RA forward pass. Therefore, this is not the throughput of the complete proposed segmentation pipeline. The full end-to-end latency should be reported.\end{revcomment}
\textbf{Response.} We agree; the original 10 images/s excluded the LU-Net+RA forward pass. We now report a complete-pipeline benchmark on an NVIDIA GeForce RTX 4070 Ti SUPER (FP32, batch size one, 50 updates; eight images, five warm-ups and ten timed repetitions per image and mode in rotated order, CUDA-synchronised). Mean latency is 4.51~ms for LU-Net+RA, 117.95~ms for refinement, 119.02~ms for the complete model and 140.50~ms from a decoded RGB image to a binary mask in CPU memory (7.12 images/s in serial execution); the last figure includes preprocessing and transfers but excludes disk access and image decoding. A separate CUDA-event session gives 112.02~ms for the complete model. Because the stages were timed separately, their sum (122.46~ms) need not equal the directly measured complete-model time.

\textbf{Changes in the manuscript.} Section~4.9 (pages~17--18, lines~570--585); Table~14 (page~17).

\subsection*{Comment 13}
\begin{revcomment}The repeated-sampling experiment is quantitatively underdeveloped. Variance is shown for one CHASE-DB1 image, but the manuscript does not provide the number of samples, nor a quantitative relationship between variance and segmentation error. Reporting the variance--error correlation over the complete test set would make this analysis substantially more meaningful.\end{revcomment}
\textbf{Response.} We removed the single-image variance figure and all calibration statements, because the sample count and display scale of that figure could not be documented. We did not perform a variance--error analysis, and the manuscript makes no uncertainty or calibration claim. Repeated sampling (five noise seeds per model, and 20 in the update-count comparison) is used only to average metrics.

\textbf{Changes in the manuscript.} Section~3.3.6 (pages~9--10, lines~319--341); Section~4.8 (page~17, lines~552--568); Section~5.1 (pages~19--20, lines~632--659).

\subsection*{Comment 14}
\begin{revcomment}The absence of augmentation is potentially important given the small datasets. The model uses no image, geometric, or mask augmentation, despite training on only 20 CHASE-DB1 and 20 DRIVE images (and 33 HRF images). The authors should either justify this choice or provide an augmentation ablation.\end{revcomment}
\textbf{Response.} The coarse LU-Net+RA is trained with aligned rotations and flips, whereas the mask autoencoder and the diffusion refiner are trained without augmentation. We did not run an augmentation ablation for the refiner, so we cannot quantify its effect or claim that augmentation is unnecessary; this is stated in the limitations. The spatial residual refiner was trained with aligned flips and right-angle rotations, so its comparison with the encoded-grid refiners also differs in augmentation, which we now note.

\textbf{Changes in the manuscript.} Section~4.2 (page~11, lines~367--393); Section~4.6 (page~16, lines~516--537); Section~5.1 (pages~19--20, lines~632--659).

\section*{Reviewer 2}
\subsection*{Comment 1}
\begin{revcomment}The most critical issue is the absence of paired comparison. The authors use published LU-Net+RA results as context, but they do not evaluate the coarse prediction and the refined prediction on the same test images under the same protocol. Without paired results, one cannot tell whether the Dice improvement comes from the diffusion refiner or from differences in dataset splits and evaluation procedures. The authors must add coarse metrics on the same split and report per-image Dice changes, means, and confidence intervals.\end{revcomment}
\textbf{Response.} Coarse and refined predictions are now evaluated on the same test images with identical annotations, preprocessing, FOV masks, threshold and metric code. The saved global/frozen records underlying Table~6 provide the paired analysis; no scores from a different refiner variant are substituted. With five training seeds and five sampling seeds, mean Dice is 0.811, 0.792 and 0.732 for LU-Net+RA and 0.746, 0.790 and 0.666 after diffusion refinement on CHASE-DB1, DRIVE and HRF. Table~6 now includes exploratory paired crossed-bootstrap 95\% intervals, and Tables~7--9 provide all 40 per-image means and paired differences. All three mean changes and their exploratory intervals are negative. Thus the controlled comparison does not support a refinement benefit. Section~4.4.1 gives the aggregation, bootstrap procedure and small-sample limitations; the historical positive scores and their tests are not transferred to this comparison.

\textbf{Changes in the manuscript.} Section~4.1 (pages~10--11, lines~344--365); Section~4.2 (page~11, lines~367--393); Section~4.3 (pages~11--12, lines~395--416); Section~4.4 (page~12, lines~418--432); Section~4.4.1 (pages~12--13, lines~434--453); Table~6 (page~12); Table~7 (page~13); Table~8 (page~13); Table~9 (page~14).

\subsection*{Comment 2}
\begin{revcomment}All results are single-run point estimates. There are no repeated seeds, no standard deviations, and no statistical tests. The Dice change from 10 to 100 DDIM transitions is small, and a single run cannot establish a reliable trend. At least five repeated runs with different random seeds are needed, with means, standard deviations, and paired tests for key comparisons.\end{revcomment}
\textbf{Response.} The reference and matched encoded-grid configurations use five separate training runs (seeds 42, 123, 2026, 3000 and 3001), and stochastic modes use five matched initial-noise seeds per trained model. We report means and sample SDs across the five training-seed means. Component, spatial-refiner and OOF comparisons report paired tests with the stated Holm families; Table~6 and Section~4.4.1 provide exploratory paired crossed-bootstrap intervals for the principal coarse-versus-refiner comparison. No new paired hypothesis tests were added to that principal comparison. With five nonzero untied pairs, the minimum two-sided exact signed-rank probability is 0.0625. Parametric tests require the corresponding independence and distributional assumptions and are not substitutes for stronger replication. The separate transition-count experiment uses one fixed trained model and twenty noise seeds: it shows negligible observed Dice change from 50 to 100 updates, with an interval including zero, rather than equivalence or a five-training-run trend. Counts below 50 were not evaluated in that protocol.

\textbf{Changes in the manuscript.} Section~4.2 (page~11, lines~367--393); Section~4.3 (pages~11--12, lines~395--416); Section~4.4.1 (pages~12--13, lines~434--453); Section~4.5 (pages~15--16, lines~481--514); Section~4.4.2 (pages~14--15, lines~455--479); Section~4.8 (page~17, lines~552--568).

\subsection*{Comment 3}
\begin{revcomment}There is no controlled comparison with existing diffusion refinement methods. The paper mentions HiDiffSeg, RefineCatDiff, and the plug-and-play diffusion denoising module of Li et al., but does not compare against them under the same setting. If reproduction is not possible, the authors should at least explain why and include a deterministic learned refiner as a baseline to isolate the benefit of the diffusion formulation.\end{revcomment}
\textbf{Response.} We added two deterministic learned refiners, as the reviewer suggested: a deterministic refiner on the same encoded grid and a spatial residual U-Net that operates on the RGB image and the coarse map. The deterministic encoded-grid refiner outperforms the diffusion refiner on all three datasets (by 0.033, 0.005 and 0.063 Dice), so we find no benefit of the diffusion formulation in this setting. We did not reproduce HiDiffSeg, RefineCatDiff or the module of Li et al.: they diffuse different quantities (skeleton-guided hierarchical maps, categorical labels and intermediate feature maps), so each would need a separate implementation and tuning. Because our refiner does not improve its own coarse input, we make no claim relative to these methods, and we state this in Section~2.

\textbf{Changes in the manuscript.} Section~2 (pages~3--4, lines~89--137); Section~4.5 (pages~15--16, lines~481--514); Table~11 (page~15); Section~4.6 (page~16, lines~516--537); Table~12 (page~16).

\subsection*{Comment 4}
\begin{revcomment}There are no component ablations. For example, there is no target-diffusion baseline, no comparison of conditioning injection methods, and no comparison between freezing and fine-tuning the mask autoencoder. Without these ablations, the design choices are not supported. I suggest at least four ablations, including a deterministic refiner, target diffusion, no image conditioning, and spatial conditioning instead of a global vector.\end{revcomment}
\textbf{Response.} We added the four suggested ablations and two more, each with five training seeds on all three datasets (105 refiners): deterministic refinement, target diffusion (from the coarse map and from Gaussian noise), no image conditioning, spatial FiLM conditioning, additive instead of FiLM injection, and joint fine-tuning of the mask autoencoder. Conditioning variants change Dice by less than 0.0002. Target diffusion from the coarse start gives similar mean Dice to the reference and fails from a Gaussian start (Dice 0.03--0.06). The deterministic refiner and autoencoder fine-tuning improve on the diffusion reference but exceed the coarse baseline only on DRIVE. None of the variants supports the original design choices as optimal.

\textbf{Changes in the manuscript.} Section~4.5 (pages~15--16, lines~481--514); Table~11 (page~15); Section~5 (page~19, lines~599--630).

\subsection*{Comment 5}
\begin{revcomment}The dataset splits and evaluation protocol are unclear. Is the CHASE-DB1 20/8 split subject-level independent. Is the HRF 33/12 split stratified by pathology. There is no validation set and no checkpoint selection strategy. The evaluation threshold, the FOV mask provenance, and whether results are averaged per image or computed from pooled pixels are all unknown. These details must be provided.\end{revcomment}
\textbf{Response.} The partitions use split seed 42 and are fixed across training seeds. CHASE-DB1 is split by subject into 16/4/8 fitting/validation/test images (8/2/4 subjects), keeping both eyes together; DRIVE keeps its official 20-image test set and splits the training pool 16/4; HRF uses 27/6/12 images, stratified 9/2/4 within each of the healthy, diabetic-retinopathy and glaucoma groups. The coarse model is selected by validation mean Dice (at most 1000 epochs, patience 20); the reference and matched encoded-grid refiners use the predefined final epoch 100, and no test result is used for selection. Predictions are binarised at 0.5, and metrics are computed per image and averaged rather than pooled. CHASE-DB1 is evaluated on the whole image with first-observer annotations; DRIVE uses the dataset FOV masks; HRF uses the FOV bitmap in the Supervisely annotation files, whose conversion history is not documented. These partitions define the reference fitting, validation and test sets. The OOF study further subdivides the fitting set into folds while retaining the fixed test set; the spatial refiner has a separate validation-based selection recipe.

\textbf{Changes in the manuscript.} Section~4.1 (pages~10--11, lines~344--365); Table~5 (page~10); Section~4.2 (page~11, lines~367--393); Section~4.3 (pages~11--12, lines~395--416); Section~4.4.2 (pages~14--15, lines~455--479).

\subsection*{Comment 6}
\begin{revcomment}The coarse predictions used to train the refiner are in-sample predictions from images that LU-Net+RA was trained on. This creates a train-test shift. The refiner should be trained on out-of-fold coarse predictions, or at least the impact should be discussed and tested.\end{revcomment}
\textbf{Response.} We added the completed matched OOF experiment to Section~4.4.2 and Table~10. It comprises 55 fold-baseline fits and 15 OOF-refiner fits: four subject-grouped folds for CHASE-DB1, four image folds for DRIVE, and three pathology-stratified folds for HRF, with fold membership fixed using seed 20261009. Fold predictors exclude their held-out images from fitting and use the separate original validation partition for checkpoint selection. The in-sample arm is the global/frozen reference of Table~6; both arms share fixed-test coarse predictions, the initial frozen autoencoder, refiner initialisation and matched optimisation and sampling settings. In-sample and OOF Dice means are 0.746145777 and 0.745915967 on CHASE-DB1, 0.789757278 and 0.789614005 on DRIVE, and 0.665725620 and 0.665522152 on HRF. OOF-minus-in-sample intervals are $[-0.000524367,-0.000031774]$, $[-0.000286638,-0.000032637]$, and $[-0.000681909,+0.000232154]$, respectively. None of the Holm-corrected training-seed tests reaches 0.05. These existing OOF calculations were incorporated without recalculation. Nonsignificance does not establish equivalence or absence of input shift: reduced fold-fitting size, frozen DRIVE autoencoder pretraining on in-sample coarse maps, small test cohorts and limited training replication remain limitations.

\textbf{Changes in the manuscript.} Section~4.2 (page~11, lines~367--393); Section~4.4.2 (pages~14--15, lines~455--479); Table~10 (page~14); Section~5.1 (pages~19--20, lines~632--659).

\subsection*{Comment 7}
\begin{revcomment}The term feature space is misleading. The encoded grid is spatially downsampled by a factor of four, but it contains about four times more scalar values than the input mask, so it is not a compressed representation. The authors should describe it as a spatially downsampled multi-channel representation and avoid terms like latent or compressed.\end{revcomment}
\textbf{Response.} We now use ``spatially downsampled multi-channel representation'' throughout the text and in Figure~1, and no longer describe the grid as latent, compressed or a feature space.

\textbf{Changes in the manuscript.} Section~1 (pages~1--2, lines~20--78); Section~3.3 (page~5, lines~180--193); Section~3.3.3 (page~7, beginning at line~254); Figure~1 (page~3).

\subsection*{Comment 8}
\begin{revcomment}The repeated-sampling variability analysis is only qualitative. The number of samples, display scale, and calibration metrics are missing. If this analysis is kept, it should be made quantitative, for example by relating variance to error or by providing reliability diagrams. Otherwise, the calibration implications should be removed.\end{revcomment}
\textbf{Response.} We removed the qualitative variance figure and the calibration statements from the abstract, introduction, contributions and conclusion. We did not compute reliability diagrams or a variance--error relationship, and the manuscript makes no calibration claim.

\textbf{Changes in the manuscript.} Section~3.3.6 (pages~9--10, lines~319--341); Section~5.1 (pages~19--20, lines~632--659).

\subsection*{Comment 9}
\begin{revcomment}Many implementation details are missing, including the refinement strength rho, the evaluation seed, interpolation kernels, batch size, weight decay, warm-up, FOV mask provenance, and the layer-by-layer diffusion U-Net specification. The authors acknowledge this in the limitations, but the revised manuscript should provide as much as possible and release code and weights.\end{revcomment}
\textbf{Response.} The reference uses start index 99 of a 1000-step cosine schedule, 50 DDIM updates, $\eta=0$ and guidance 1.5. Target diffusion from pure Gaussian noise instead begins at index 999; the deterministic control uses no DDIM. The strength $\rho$ of the earlier run was not recorded. The revised manuscript documents the training and sampling seeds, bilinear interpolation for images and nearest-neighbour interpolation for masks, batch size four, the AdamW learning rates and weight decay, cosine scheduling, gradient clipping, conditioning dropout, checkpoint selection, the FOV sources and a layer-level denoiser specification. Code and evaluation scripts are public at \url{https://github.com/thebabak/Refining_Retinal_Vessel_Segmentation_with_Conditional_Diffusion_Models}, and retained model weights and saved repeated-training predictions are available from the corresponding authors on reasonable request. The fixed-checkpoint repeated-sampling experiment is reported from archived score and timing records; complete probability arrays for every repeat were not retained for fresh mask-level verification.

\textbf{Changes in the manuscript.} Section~3.3.4 (page~8, lines~272--283); Table~3 (page~8); Section~3.3.5 (pages~8--9, lines~285--317); Section~3.3.6 (pages~9--10, lines~319--341); Section~4.1 (pages~10--11, lines~344--365); Table~4 (page~10); Section~4.2 (page~11, lines~367--393); Data Availability Statement (page~20, ending at line~680).

\bigskip
We hope that the revised manuscript addresses the reviewers' concerns, and we thank them again for comments that substantially improved the study.

\medskip
Sincerely,\\
Hanieh Naderi, on behalf of all authors

\end{document}
